\section{Introduction}
\label{section:introduction}

\subsection{Motivation}
\label{sec:intro:motivation}

Large Vision-Language Models (LVLMs) have shown growing capabilities in perceiving and synthesizing visual and textual information, with rapid progress across general-purpose~\cite{bai2025qwen3, achiam2023gpt, singh2025openai, liu2024deepseek} and medical~\cite{li2024llavamed, chen2024chexagent, chen2024towards, sellergren2025medgemma, alsaad2024multimodal, saab2024capabilities} domains. A model that can read a chest radiograph and produce a structured report, or answer a clinician's question about a CT slice, is an obvious candidate for relieving pressure on radiology workflows and for extending access to imaging expertise in settings where a specialist is not available.

In high-stakes settings such as healthcare, however, LVLMs continue to exhibit unintended failure modes. They hallucinate clinical findings that are not present in the image, produce responses that are linguistically fluent yet factually unsupported by the underlying visual evidence~\cite{zhang2024medihall, seth2025hallucinogen}, and fail to reliably follow clinician instructions~\cite{vadlapati2026ai, xia2024cares}. These failures undermine medical reliability and limit safe deployment in real-world clinical workflows.

What makes these failures difficult is not their frequency but their form. Medical pretraining is highly effective at teaching a model the vocabulary, register, and structure of clinical text, so a hallucinated finding arrives phrased exactly as a radiologist would phrase it, hedged in the appropriate places and embedded in an otherwise accurate description. The error is confined to the few words that carry the diagnosis. This is a worse failure mode than incoherence would be, because fluency and confidence are the cues a human reader ordinarily uses to gauge reliability, and here they have been decoupled from correctness.

A related and equally serious problem is invisible at the level of the output altogether: a model may produce the correct answer for reasons that have nothing to do with the image, relying on the base rate of a finding or the anatomical default for a modality. Such an answer is indistinguishable from a well-grounded one when read, but it will fail on the atypical case --- which is precisely the case where a clinician most needs help. This failure is also self-concealing under standard evaluation: a model exploiting the base rates of its training distribution will score well on a test set drawn from the same distribution, and the deficit surfaces only on the rare presentations that aggregate accuracy barely registers.

Both problems share a structure worth naming at the outset. The model's output distribution is well-formed but incorrectly conditioned: it has learned what a good answer looks like without having learned what makes a superficially identical answer wrong. Supervised fine-tuning cannot easily repair this, because a demonstration exhibits only the correct answer and never the near-miss that must be rejected. Preference-based post-training is attractive precisely because a preference pair carries the information a demonstration lacks --- a negative example close enough to the positive one to isolate the distinction that matters.

\begin{figure}[t]
    \centering
    \includegraphics[width=0.75\textwidth]{generations.pdf}
    \caption{Reward assignment and visual attention in DPO versus \ourloss{}. DPO exhibits the same diffused visual attention as the base LVLM and fails to prioritize the reward of the clinically decisive words ``left lung'' (darker is higher).}
    \label{fig:framework_overview}
\end{figure}

\subsection{Preference Optimization and Its Clinical Mismatch}
\label{sec:intro:mismatch}

Beyond the initial supervised learning phase, post-training techniques have demonstrated remarkable success in refining model behaviour and improving task-specific performance~\cite{rafailov2023direct, ethayarajh2023kto, azar2024general}. Among them, Direct Preference Optimization (DPO) has gained popularity for its ability to skip explicit reward modeling~\cite{hein-etal-2025-chexalign, zhu2024mmedpo, savage2025fine, sun2024self}. Because it requires only comparisons between outputs rather than gold-standard demonstrations, it appears well suited to a domain where writing a reference answer requires a specialist but ranking two candidates does not.

This thesis argues that standard DPO and its variants face three critical limitations when applied to the medical domain, and that these limitations are not incidental but follow from the structure of the objective.

\emph{First}, the pairwise preference format is too coarse to capture the feedback clinicians actually provide, causing DPO-based approaches to yield inconsistent gains over supervised fine-tuning~\cite{kim2026benchmarking}. While clinical evaluation hinges on fine-grained correctness --- exact terminology, accurate anatomical localization --- standard DPO treats all response tokens uniformly during optimization. Consequently, diagnostically important spans such as findings, anatomical locations, and measurements contribute no more to the training signal than generic filler phrases such as ``the image shows'' or ``there is evidence of''. Figure~\ref{fig:framework_overview} shows this directly: the reward a DPO-aligned model assigns is spread evenly across the response, with the decisive words receiving no special weight.

\emph{Second}, because AI-generated preference labels are unreliable in medicine~\cite{delucia2026same, rankin2020reliability}, prior work substitutes the supervised fine-tuning ground truth as the preferred response~\cite{zhu2024mmedpo, liu2025rrg, savage2025fine}. This creates an off-policy distribution shift~\cite{yang2025mitigating}. Clinical reference answers are short and structured; model outputs are long and explanatory. The systematic stylistic gap between them becomes a spurious signal that the model exploits as a reward-hacking shortcut, optimizing stylistic patterns rather than clinical correctness. Because the preference objective is defined purely by the contrast between two responses, any property that separates them reliably will be learned --- and imitating a register is far easier than learning a diagnosis.

\emph{Third}, existing alignment objectives lack explicit visual grounding constraints, leaving models insensitive to subtle yet diagnostically decisive pathological features. Nothing in the objective distinguishes a correct answer read from the image from an identical answer produced by a linguistic prior. As Figure~\ref{fig:framework_overview} illustrates, DPO exhibits the same diffused visual attention as the unaligned base model even where it has corrected the answer.

Taken together these three limitations suggest that the inconsistent benefit of preference optimization in medicine is not a matter of insufficient tuning or insufficient data. Each is a consequence of a design decision that is sound in the general domain and unsound here, and each admits a specific remedy. That is the position this thesis develops and tests.

\subsection{Scope}
\label{sec:intro:scope}

Three boundaries are worth stating before the argument begins.

This thesis concerns \emph{post-training alignment}, not pretraining or architecture. The visual encoder and cross-modal projector of both backbones remain frozen throughout, and only the language model is adapted, using low-rank updates. One consequence is methodological rather than incidental: because the perceptual front end never changes, any improvement in visual grounding reported in Chapter~\ref{section:results} must be attributable to how the model \emph{uses} the visual representation rather than to a better representation, which makes the grounding claim cleaner than it would otherwise be. The corresponding cost is that errors rooted in perception rather than in verbalization are outside the reach of the method entirely.

The evaluation is confined to radiology --- two visual question answering datasets and one report generation dataset --- and is in-distribution, on the held-out test split of the corpus each model was aligned on. Whether the framework transfers to pathology, dermatology, or ophthalmology, whose visual grounding characteristics differ substantially, is left open, and is revisited in Section~\ref{sec:con:limitations}.

Finally, this work aims at factual and evidential correctness, not at clinical deployment. Nothing here addresses calibration, uncertainty communication, regulatory validation, or the human factors of how a clinician should incorporate model output into a decision. Improving factuality is a precondition for those questions rather than an answer to them, and Section~\ref{sec:con:impact} returns to the risk that improving surface quality alongside factuality can make residual errors harder, not easier, for a reader to catch.

\subsection{Thesis Statement and Research Questions}
\label{sec:intro:rqs}

\noindent\textbf{Thesis statement.} \emph{The inconsistent benefit of preference optimization for medical vision-language models is not a limitation of preference learning as such, but a consequence of three specific mismatches --- reward granularity, preference-pair construction, and the absence of a grounding constraint. Addressing all three jointly, within a single objective, yields consistent improvements in clinical factuality and in measured visual grounding across backbones and tasks.}

The thesis is organized around four research questions.

\begin{itemize}
    \item \textbf{RQ1.} Does the sequence-level reward of standard DPO fail to concentrate learning signal on the clinically decisive spans of a response, and does restricting the reward to those spans improve clinical correctness? \emph{(Sections~\ref{sec:ana:coarse}, \ref{sec:met:ranking}, \ref{sec:res:main}.)}

    \item \textbf{RQ2.} When the supervised reference is used as the preferred response, does the model optimize response style in preference to clinical correctness, and does constructing preference pairs from the model's own outputs remove this confound? \emph{(Sections~\ref{sec:ana:reward-hacking}, \ref{sec:met:data}, \ref{sec:res:dynamics}, \ref{sec:res:abl-style}.)}

    \item \textbf{RQ3.} How does the direction of the KL regularizer govern the trade-off between clinical precision and generative diversity, and is a bidirectional regularizer superior to either direction alone? \emph{(Sections~\ref{sec:ana:kl}, \ref{sec:met:kl}, \ref{sec:res:abl-kl}.)}

    \item \textbf{RQ4.} Can an explicit visual-contrastive objective, built on lesion-targeted rather than global image corruption, improve a model's reliance on diagnostically decisive image regions as measured inside the network? \emph{(Sections~\ref{sec:ana:grounding}, \ref{sec:met:corruption}, \ref{sec:res:grounding}, \ref{sec:res:abl-visual}.)}
\end{itemize}

\subsection{Contributions}
\label{sec:intro:contributions}

\begin{itemize}
    \item \textbf{An analysis of why preference optimization underperforms in clinical settings.} We characterize three failure modes --- coarse reward assignment, stylistic reward hacking, and the absence of grounding constraints --- and add a fourth axis of analysis in the geometry of the KL regularizer. We show that stylistic reward hacking is visible not only in downstream accuracy but in the optimization itself: an off-policy DPO run saturates its preference objective within a twelfth of an epoch, after which its gradients collapse by six orders of magnitude and training effectively halts.

    \item \textbf{A medical preference data framework.} We propose a strategy for generating on-policy preference datasets from existing medical SFT data. By creating counterfactual pairs that correct specific model hallucinations while preserving original syntax, we move beyond simple ground-truth comparisons to provide more precise supervision. The construction additionally produces lesion-targeted corrupted images by text-guided medical segmentation, giving a visual counterpart to each textual pair.

    \item \textbf{A multimodal fine-grained regularized loss.} We introduce \ourloss{}, a loss operating at both the text and image levels. It provides fine-grained medical rewards on clinically divergent spans, applies a bidirectional token-wise KL regularizer that balances coverage preservation against tail suppression, and adds a dual-constraint visual term that reverses the preference direction when the supporting evidence has been destroyed.

    \item \textbf{Empirical validation across backbones and tasks.} Experiments across medical VQA and report-generation benchmarks demonstrate that \ourloss{} consistently outperforms both standard DPO and RRPO, achieving an average relative improvement of $10.24\%$ across two state-of-the-art LVLMs, together with layer-wise evidence of improved visual grounding and controlled ablations isolating each component.
\end{itemize}

\subsection{Thesis Outline}
\label{sec:intro:outline}

Chapter~\ref{section:background} provides the technical background: the architecture of vision-language models, their medical specialization, the derivation of DPO from reinforcement learning from human feedback, the information-theoretic properties of the KL divergence in both directions, and the benchmarks and metrics used to evaluate clinical systems. Chapter~\ref{section:related} positions this work against the literature on fine-grained alignment, multimodal preference learning, and clinical preference optimization, closing with a comparison table.

Chapter~\ref{section:analysis} develops the analysis that motivates the method, treating each failure mode in turn and concluding with a set of design desiderata. Chapter~\ref{section:method} presents \ourloss{}: the on-policy preference construction pipeline, the lesion-targeted visual corruption procedure, the dual-constraint ranking loss, the bidirectional token-wise KL regularizer, and the implementation.

Chapter~\ref{section:setup} describes the experimental setup, including a characterization of the generated preference corpora. Chapter~\ref{section:results} reports the main results, the layer-wise grounding analysis, the training dynamics, and four ablations isolating the individual design choices. Chapter~\ref{section:conclusion} discusses what the results support, states the limitations of the study, addresses broader impact, and outlines future work.
